\documentclass[11pt,a4paper]{article}

\usepackage[utf8]{inputenc}
\usepackage[T1]{fontenc}
\usepackage{lmodern}
\usepackage[margin=1in]{geometry}
\usepackage{amsmath,amssymb,amsthm}
\usepackage{mathtools}
\usepackage{bm}
\usepackage{hyperref}
\usepackage{enumitem}
\usepackage{booktabs}
\usepackage{tcolorbox}
\usepackage{fancyhdr}

\hypersetup{
    colorlinks=true,
    linkcolor=blue!70!black,
    urlcolor=blue!60!black
}

\theoremstyle{definition}
\newtheorem{definition}{Definition}[section]
\newtheorem{theorem}{Theorem}[section]
\newtheorem{proof_}{Proof}[section]
\newtheorem{remark}{Remark}[section]

\newtcolorbox{intuition}{
    colback=blue!5!white,
    colframe=blue!50!black,
    title=Geometric Intuition,
    fonttitle=\bfseries
}

\newtcolorbox{application}{
    colback=green!5!white,
    colframe=green!50!black,
    title=Application to FB-CycleGAN,
    fonttitle=\bfseries
}

\pagestyle{fancy}
\fancyhf{}
\fancyhead[L]{\small Fourier Analysis Reference}
\fancyhead[R]{\small CS424 FB-CycleGAN}
\fancyfoot[C]{\thepage}

\title{\textbf{Fourier Analysis from First Principles} \\[0.5em]
\Large A Self-Contained Reference for Signal and Image Analysis}
\author{CS424 Group Project --- Internal Reference}
\date{March 2026}

\begin{document}
\maketitle
\tableofcontents
\newpage

% ======================================================================
\section{Complex Numbers as 2D Objects}
% ======================================================================

Before Fourier analysis, we need to be precise about complex numbers.

\begin{definition}[Complex Number]
A complex number $z = a + bi$ where $a, b \in \mathbb{R}$ and $i = \sqrt{-1}$ is a point in the 2D \textbf{complex plane}:
\begin{itemize}
    \item The \textbf{real part} $\mathrm{Re}(z) = a$ is the $x$-coordinate.
    \item The \textbf{imaginary part} $\mathrm{Im}(z) = b$ is the $y$-coordinate.
    \item The \textbf{modulus} $|z| = \sqrt{a^2 + b^2}$ is the distance from the origin.
    \item The \textbf{conjugate} $\overline{z} = a - bi$ reflects across the real axis.
\end{itemize}
\end{definition}

\begin{definition}[Euler's Formula]
\begin{equation}
    e^{i\theta} = \cos\theta + i\sin\theta
\end{equation}
This maps the angle $\theta$ to a point on the unit circle in the complex plane. As $\theta$ increases from $0$ to $2\pi$, $e^{i\theta}$ traces the full unit circle:
\begin{align}
    \theta = 0 &\implies e^{i \cdot 0} = 1 + 0i = (1, 0) \quad\text{(right)} \\
    \theta = \pi/2 &\implies e^{i\pi/2} = 0 + 1 \cdot i = (0, 1) \quad\text{(top)} \\
    \theta = \pi &\implies e^{i\pi} = -1 + 0i = (-1, 0) \quad\text{(left)} \\
    \theta = 3\pi/2 &\implies e^{i3\pi/2} = 0 - 1 \cdot i = (0, -1) \quad\text{(bottom)}
\end{align}
\end{definition}

\begin{definition}[Complex Exponential at Frequency $k$]
The function $e^{ikt}$ traces the unit circle $k$ times as $t$ goes from $0$ to $2\pi$:
\begin{itemize}
    \item $k = 1$: one full loop (slow spin)
    \item $k = 3$: three full loops (fast spin)
    \item $k = -2$: two loops in the opposite direction
\end{itemize}
Its real part is $\cos(kt)$ and its imaginary part is $\sin(kt)$, so $e^{ikt}$ packages a cosine and sine at frequency $k$ into one object.
\end{definition}

% ======================================================================
\section{Function Spaces and Inner Products}
% ======================================================================

\begin{definition}[$L^2$ Space]
$L^2([0, 2\pi])$ is the set of all square-integrable functions on $[0, 2\pi]$:
\begin{equation}
    L^2([0, 2\pi]) = \left\{ f : [0, 2\pi] \to \mathbb{C} \;\middle|\; \int_0^{2\pi} |f(t)|^2 \, dt < \infty \right\}
\end{equation}
This is a vector space: you can add functions and scale them, and the results are still in $L^2$.
\end{definition}

\begin{definition}[Inner Product on $L^2$]
For $f, g \in L^2([0, 2\pi])$:
\begin{equation}
    \langle f, g \rangle = \frac{1}{2\pi} \int_0^{2\pi} f(t) \, \overline{g(t)} \, dt
\end{equation}
The $\frac{1}{2\pi}$ is a normalisation constant so that $\langle e^{ikt}, e^{ikt} \rangle = 1$.
\end{definition}

\begin{remark}[Why Divide by $2\pi$?]
Without it, $\langle e^{ikt}, e^{ikt} \rangle = \int_0^{2\pi} |e^{ikt}|^2 \, dt = \int_0^{2\pi} 1 \, dt = 2\pi$. We divide by $2\pi$ to make the basis vectors unit length. This is a normalisation convention---some texts put $\frac{1}{\sqrt{2\pi}}$ on both the forward and inverse transforms instead. The choice doesn't matter as long as it's consistent.
\end{remark}

\subsection{Inner Product as Projection}

In $\mathbb{R}^2$, the dot product $\langle \mathbf{v}, \mathbf{u} \rangle$ where $\mathbf{u}$ is a unit vector gives the length of $\mathbf{v}$'s shadow (projection) onto $\mathbf{u}$:
\begin{equation}
    \mathbf{v} = (3, 4), \quad \mathbf{e}_1 = (1, 0) \quad\implies\quad \langle \mathbf{v}, \mathbf{e}_1 \rangle = 3 \cdot 1 + 4 \cdot 0 = 3
\end{equation}
The result $3$ is the component of $\mathbf{v}$ along $\mathbf{e}_1$.

The $L^2$ inner product does exactly the same thing for functions. $\langle f, e^{ikt} \rangle$ measures how much of $f$ ``points in the direction'' of $e^{ikt}$---i.e., how much of frequency $k$ is present in $f$.

% ======================================================================
\section{Orthonormality of the Exponential Basis}
% ======================================================================

\begin{theorem}[Orthonormality]
The functions $\{e_k(t) = e^{ikt}\}_{k \in \mathbb{Z}}$ are orthonormal in $L^2([0, 2\pi])$:
\begin{equation}
    \langle e_k, e_m \rangle = \begin{cases} 1 & \text{if } k = m \\ 0 & \text{if } k \neq m \end{cases}
\end{equation}
\end{theorem}

\textbf{Proof.} Compute directly:
\begin{equation}
    \langle e_k, e_m \rangle = \frac{1}{2\pi} \int_0^{2\pi} e^{ikt} \cdot \overline{e^{imt}} \, dt = \frac{1}{2\pi} \int_0^{2\pi} e^{ikt} \cdot e^{-imt} \, dt = \frac{1}{2\pi} \int_0^{2\pi} e^{i(k-m)t} \, dt
\end{equation}

\textbf{Case $k = m$:} The exponent is $0$, so $e^{i \cdot 0 \cdot t} = 1$:
\begin{equation}
    \frac{1}{2\pi} \int_0^{2\pi} 1 \, dt = \frac{2\pi}{2\pi} = 1 \quad\checkmark
\end{equation}

\textbf{Case $k \neq m$:} Let $p = k - m \neq 0$. By Euler's formula, $e^{ipt} = \cos(pt) + i\sin(pt)$:
\begin{align}
    \int_0^{2\pi} e^{ipt} \, dt &= \int_0^{2\pi} \cos(pt) \, dt + i \int_0^{2\pi} \sin(pt) \, dt
\end{align}
Since $p$ is a nonzero integer, both $\cos(pt)$ and $\sin(pt)$ complete full cycles on $[0, 2\pi]$. The integral of any full cycle of sine or cosine is zero (the positive half cancels the negative half):
\begin{align}
    \int_0^{2\pi} \cos(pt) \, dt &= \left[\frac{\sin(pt)}{p}\right]_0^{2\pi} = \frac{\sin(2\pi p) - \sin(0)}{p} = 0 \\
    \int_0^{2\pi} \sin(pt) \, dt &= \left[\frac{-\cos(pt)}{p}\right]_0^{2\pi} = \frac{-\cos(2\pi p) + \cos(0)}{p} = \frac{-1+1}{p} = 0
\end{align}
Therefore $\langle e_k, e_m \rangle = 0$ when $k \neq m$. $\square$

\begin{intuition}
Two sinusoids at different integer frequencies always cancel over a full period. If you multiply $\cos(2t)$ by $\cos(3t)$ and integrate, the product oscillates symmetrically around zero. Only when the frequencies match does the product $|e^{ikt}|^2 = 1$ stay positive everywhere.
\end{intuition}

% ======================================================================
\section{The Fourier Transform}
% ======================================================================

\subsection{Fourier Coefficients = Projections onto the Basis}

\begin{definition}[Fourier Coefficient]
For $f \in L^2([0, 2\pi])$, the $k$-th Fourier coefficient is the projection of $f$ onto the basis function $e_k$:
\begin{equation}
    \hat{f}[k] = \langle f, e_k \rangle = \frac{1}{2\pi} \int_0^{2\pi} f(t) \, e^{-ikt} \, dt
\end{equation}
\end{definition}

\begin{intuition}
The $e^{-ikt}$ factor ``wraps'' the signal around the origin at speed $k$ (the minus sign reverses the direction). The integral computes the centre of mass of the wrapped signal. If $f$ contains a component at frequency $k$, the peaks consistently land on one side of the circle, pulling the centre of mass away from the origin. If not, the peaks cancel, and $\hat{f}[k] \approx 0$.

This is precisely the 3Blue1Brown visualisation.
\end{intuition}

\subsection{Reconstruction (Inverse Fourier Transform)}

\begin{theorem}[Fourier Series]
Any $f \in L^2([0, 2\pi])$ can be reconstructed from its coefficients:
\begin{equation}
    f(t) = \sum_{k=-\infty}^{\infty} \hat{f}[k] \, e^{ikt}
\end{equation}
\end{theorem}

\textbf{Derivation.} Assume $f$ can be expanded as $f(t) = \sum_k c_k \, e^{ikt}$. To find $c_m$, take the inner product of both sides with $e_m$:
\begin{align}
    \langle f, e_m \rangle &= \left\langle \sum_k c_k \, e_k, \; e_m \right\rangle \\
    &= \sum_k c_k \, \langle e_k, e_m \rangle \quad\text{(linearity of inner product)} \\
    &= c_m \cdot 1 + \sum_{k \neq m} c_k \cdot 0 \quad\text{(orthonormality)} \\
    &= c_m
\end{align}
So $c_m = \langle f, e_m \rangle = \hat{f}[m]$. The forward transform extracts coefficients; the inverse sums them back. $\square$

\begin{remark}
This is identical to how in $\mathbb{R}^3$, if $\mathbf{v} = a\mathbf{e}_1 + b\mathbf{e}_2 + c\mathbf{e}_3$, then $a = \langle \mathbf{v}, \mathbf{e}_1 \rangle$. The only difference is that $L^2$ is infinite-dimensional.
\end{remark}

% ======================================================================
\section{The Discrete Fourier Transform (DFT)}
% ======================================================================

For a finite signal of $N$ samples (e.g., $N = 256$ pixels), the function space is $\mathbb{C}^N$---an ordinary finite-dimensional vector space.

\begin{definition}[DFT]
For a signal $f = (f[0], f[1], \ldots, f[N-1]) \in \mathbb{C}^N$:
\begin{equation}
    \hat{f}[k] = \sum_{n=0}^{N-1} f[n] \cdot e^{-i \, 2\pi \, kn / N}, \qquad k = 0, 1, \ldots, N-1
\end{equation}
\end{definition}

\begin{definition}[Inverse DFT]
\begin{equation}
    f[n] = \frac{1}{N} \sum_{k=0}^{N-1} \hat{f}[k] \cdot e^{+i \, 2\pi \, kn / N}, \qquad n = 0, 1, \ldots, N-1
\end{equation}
\end{definition}

\begin{remark}[DFT as Matrix Multiplication]
The DFT is a linear map $\mathbb{C}^N \to \mathbb{C}^N$ represented by the matrix $\mathbf{W}$ where:
\begin{equation}
    W_{kn} = e^{-i \, 2\pi \, kn / N}
\end{equation}
Then $\hat{\mathbf{f}} = \mathbf{W} \mathbf{f}$. The matrix $\frac{1}{\sqrt{N}} \mathbf{W}$ is unitary ($\mathbf{W}^* \mathbf{W} = N \cdot \mathbf{I}$).
\end{remark}

\subsection{Why the DFT Matrix is Unitary}

The columns of $\frac{1}{\sqrt{N}}\mathbf{W}$ are the vectors:
\begin{equation}
    \mathbf{w}_n = \frac{1}{\sqrt{N}} \left(1, \; e^{-i2\pi n/N}, \; e^{-i2\pi \cdot 2n/N}, \; \ldots, \; e^{-i2\pi(N-1)n/N}\right)
\end{equation}

Their inner product:
\begin{equation}
    \langle \mathbf{w}_n, \mathbf{w}_m \rangle = \frac{1}{N} \sum_{k=0}^{N-1} e^{-i2\pi k(n-m)/N}
\end{equation}

When $n = m$: every term is $1$, sum is $N$, result is $1$.

When $n \neq m$: this is a geometric series summing $N$-th roots of unity. The $N$-th roots of unity are evenly spaced around the unit circle and cancel to zero:
\begin{equation}
    \sum_{k=0}^{N-1} \left(e^{-i2\pi(n-m)/N}\right)^k = \frac{1 - e^{-i2\pi(n-m)}}{1 - e^{-i2\pi(n-m)/N}} = \frac{1 - 1}{1 - e^{-i2\pi(n-m)/N}} = 0
\end{equation}

Therefore the columns are orthonormal, $\mathbf{W}$ is unitary (up to scaling), and the DFT is a length-preserving change of basis. No information is lost.

\begin{remark}[Unitary $\neq$ Rotation]
A unitary matrix satisfies $\mathbf{U}^*\mathbf{U} = \mathbf{I}$, which means it preserves lengths and angles. This includes both rotations ($\det = +1$) and reflections ($\det = -1$). The precise statement is: the DFT is a \textbf{unitary change of basis} from the standard (pixel) basis to the frequency basis.
\end{remark}

% ======================================================================
\section{2D Fourier Transform (for Images)}
% ======================================================================

\begin{definition}[2D DFT]
For an image $f[m, n]$ of size $M \times N$:
\begin{equation}
    \hat{f}[u, v] = \sum_{m=0}^{M-1} \sum_{n=0}^{N-1} f[m, n] \cdot e^{-i2\pi\left(\frac{um}{M} + \frac{vn}{N}\right)}
\end{equation}
\end{definition}

\subsection{Why $e^{i(kx + ly)}$?}

The 2D basis functions are tensor products of 1D basis functions:
\begin{equation}
    e_{k,l}(x, y) = e^{ikx} \cdot e^{ily} = e^{i(kx + ly)}
\end{equation}

The last equality uses the exponential property $e^{a+b} = e^a \cdot e^b$.

\begin{itemize}
    \item $e^{ikx}$ oscillates $k$ times in the $x$ direction (constant along $y$)
    \item $e^{ily}$ oscillates $l$ times in the $y$ direction (constant along $x$)
    \item Their product oscillates in both directions simultaneously, producing stripes at angle $\arctan(l/k)$ with frequency $\sqrt{k^2 + l^2}$
\end{itemize}

This is a \textbf{tensor product} construction: if $\{e_k(x)\}$ is a basis for functions of $x$, and $\{e_l(y)\}$ is a basis for functions of $y$, then $\{e_k(x) \cdot e_l(y)\}$ is a basis for functions of $(x, y)$.

\begin{intuition}
Every image is a sum of 2D striped patterns at different frequencies, orientations, and intensities:
\begin{itemize}
    \item Low-frequency stripes ($k, l$ small): large smooth regions (brain shape, overall brightness)
    \item Mid-frequency stripes: medium structures (brain folds, organ boundaries)
    \item High-frequency stripes ($k, l$ large): fine detail (edges, texture, noise)
\end{itemize}
\end{intuition}

\subsection{Power Spectrum and Radial Average}

\begin{definition}[Power Spectrum]
\begin{equation}
    P[u,v] = |\hat{f}[u,v]|^2 = \mathrm{Re}(\hat{f}[u,v])^2 + \mathrm{Im}(\hat{f}[u,v])^2
\end{equation}
This discards the phase (position of stripes) and keeps only how much energy is at each frequency.
\end{definition}

\begin{definition}[Radial Average]
To collapse the 2D spectrum to 1D:
\begin{equation}
    P(\rho) = \frac{1}{|\{(u,v) : \sqrt{u^2+v^2} = \rho\}|} \sum_{\sqrt{u^2+v^2} = \rho} P[u,v]
\end{equation}
This averages over all orientations at frequency $\rho$, giving a single curve: spatial frequency vs.\ power.
\end{definition}

% ======================================================================
\section{Convolution}
% ======================================================================

\begin{definition}[Convolution (Continuous)]
\begin{equation}
    (f * g)(t) = \int_{-\infty}^{\infty} f(s) \cdot g(t - s) \, ds
\end{equation}
\end{definition}

\begin{definition}[Convolution (Discrete)]
\begin{equation}
    (f * g)[n] = \sum_{m=0}^{N-1} f[m] \cdot g[n - m]
\end{equation}
\end{definition}

\begin{intuition}
Convolution slides the kernel $g$ across the signal $f$. At each position $t$, it computes the overlap---the integral (or sum) of the product. The result is a weighted average where $g$ defines the weighting pattern.

Gaussian blur is convolution where $g$ is a bell curve. Each pixel becomes a weighted average of its neighbours, with nearby pixels weighted more heavily.
\end{intuition}

% ======================================================================
\section{Convolution Theorem}
% ======================================================================

\begin{theorem}[Convolution Theorem]
Convolution in the spatial domain becomes pointwise multiplication in the frequency domain:
\begin{equation}
    \widehat{f * g}[k] = \hat{f}[k] \cdot \hat{g}[k]
\end{equation}
\end{theorem}

\textbf{Proof.} Start from the Fourier transform of the convolution:
\begin{equation}
    \widehat{f * g}[k] = \int (f * g)(t) \cdot e^{-ikt} \, dt
\end{equation}

Expand the convolution:
\begin{equation}
    = \int_t \left[ \int_s f(s) \cdot g(t - s) \, ds \right] e^{-ikt} \, dt
\end{equation}

Swap the order of integration (Fubini's theorem---valid because $f, g \in L^1$ or $L^2$; for finite signals this is just swapping finite sums, which is always valid):
\begin{equation}
    = \int_s f(s) \left[ \int_t g(t - s) \cdot e^{-ikt} \, dt \right] ds
\end{equation}

For the inner integral, substitute $u = t - s$, so $t = u + s$ and $dt = du$ (since $s$ is a constant with respect to $t$, shifting by a constant doesn't change the differential):
\begin{align}
    \int_t g(t-s) \cdot e^{-ikt} \, dt &= \int_u g(u) \cdot e^{-ik(u+s)} \, du \\
    &= \int_u g(u) \cdot e^{-iku} \cdot e^{-iks} \, du \\
    &= e^{-iks} \int_u g(u) \cdot e^{-iku} \, du \\
    &= e^{-iks} \cdot \hat{g}[k]
\end{align}

The key step: $e^{-ik(u+s)} = e^{-iku} \cdot e^{-iks}$ (exponential property), and $e^{-iks}$ doesn't depend on $u$ so it factors out. The remaining integral is the Fourier transform of $g$.

Substitute back:
\begin{align}
    \widehat{f * g}[k] &= \int_s f(s) \cdot e^{-iks} \cdot \hat{g}[k] \, ds \\
    &= \hat{g}[k] \int_s f(s) \cdot e^{-iks} \, ds \\
    &= \hat{g}[k] \cdot \hat{f}[k] \qquad\square
\end{align}

\begin{remark}[Why Fubini Applies]
Fubini's theorem allows swapping the order of integration when $\int\!\!\int |h(s,t)| \, ds \, dt < \infty$. For our images, $f$ and $g$ are finite arrays of pixel values---trivially integrable. In the discrete case, we are swapping finite sums, which requires no conditions at all.
\end{remark}

% ======================================================================
\section{Eigenfunctions and Why This Basis is Special}
% ======================================================================

\begin{definition}[Eigenfunction]
For a linear operator $T$ acting on functions, $\phi$ is an eigenfunction if:
\begin{equation}
    T[\phi] = \lambda \cdot \phi
\end{equation}
where $\lambda$ is a scalar (the eigenvalue). The operator just scales $\phi$---it doesn't change its shape. This is the function-space analogue of the matrix eigenvector equation $A\mathbf{v} = \lambda\mathbf{v}$.
\end{definition}

\begin{theorem}
The complex exponentials $e^{ikt}$ are eigenfunctions of the derivative operator:
\begin{equation}
    \frac{d}{dt} e^{ikt} = ik \cdot e^{ikt}
\end{equation}
The eigenvalue is $ik$.
\end{theorem}

\begin{theorem}
The complex exponentials are eigenfunctions of \textbf{every convolution operator}. For any kernel $g$:
\begin{equation}
    (g * e^{ikt})(t) = \hat{g}[k] \cdot e^{ikt}
\end{equation}
The eigenvalue is $\hat{g}[k]$---the $k$-th Fourier coefficient of the kernel.
\end{theorem}

\begin{intuition}
In matrix algebra, if you find the eigenbasis of a matrix $A$, then $A$ becomes diagonal in that basis---each eigenvector just gets scaled by its eigenvalue. The Fourier basis is the eigenbasis of convolution. In the Fourier basis, convolution becomes diagonal: each frequency component just gets scaled by $\hat{g}[k]$. This is why the convolution theorem works.
\end{intuition}

% ======================================================================
\section{Complete Property List}
% ======================================================================

Let $\hat{f}$ denote the Fourier transform of $f$, and $\hat{f}[k]$ its $k$-th coefficient.

\subsection{Linearity}
\begin{equation}
    \widehat{af + bg}[k] = a\hat{f}[k] + b\hat{g}[k]
\end{equation}
\textit{Proof:} The integral is linear. \quad
\textit{Intuition:} Mixing two sounds gives mixed frequency content. Nothing interacts.

\subsection{Shift (Time Delay)}
\begin{equation}
    \widehat{f(t - d)}[k] = e^{-ikd} \cdot \hat{f}[k]
\end{equation}
\textit{Proof:} Substitute $u = t - d$, factor out $e^{-ikd}$. \quad
\textit{Intuition:} Shifting doesn't change \emph{what} frequencies are present ($|\hat{f}[k]|$ unchanged), only \emph{where} the waves are positioned (phase changes).

\subsection{Modulation (Frequency Shift)}
\begin{equation}
    \widehat{e^{imt} f(t)}[k] = \hat{f}[k - m]
\end{equation}
\textit{Proof:} The exponentials combine: $e^{imt} \cdot e^{-ikt} = e^{-i(k-m)t}$. \quad
\textit{Intuition:} Multiplying by a complex exponential shifts all frequencies by $m$.

\subsection{Convolution Theorem}
\begin{equation}
    \widehat{f * g}[k] = \hat{f}[k] \cdot \hat{g}[k]
\end{equation}
\textit{Proof:} Section 7. \quad
\textit{Intuition:} Filtering in space = scaling in frequency. A low-pass filter multiplies high-$k$ coefficients by $\approx 0$.

\subsection{Multiplication Theorem}
\begin{equation}
    \widehat{f \cdot g}[k] = (\hat{f} * \hat{g})[k]
\end{equation}
\textit{Intuition:} Pointwise multiplication in space = convolution in frequency. Windowing a signal smears its spectrum.

\subsection{Parseval's Theorem (Energy Conservation)}
\begin{equation}
    \int |f(t)|^2 \, dt = \sum_k |\hat{f}[k]|^2
\end{equation}
\textit{Proof:} Expand $|f|^2$ using the Fourier series, use orthonormality---cross terms vanish. \quad
\textit{Intuition:} Total energy is the same in both representations. The Fourier transform is a unitary change of basis---unitary maps preserve lengths.

\subsection{Differentiation}
\begin{equation}
    \widehat{f'}[k] = ik \cdot \hat{f}[k]
\end{equation}
\textit{Proof:} Differentiate the reconstruction $f(t) = \sum \hat{f}[k] e^{ikt}$ term by term. \quad
\textit{Intuition:} Differentiation amplifies high frequencies. Rapidly oscillating signals have large derivatives.

\subsection{Conjugate Symmetry (Real Signals)}
If $f(t)$ is real-valued:
\begin{equation}
    \hat{f}[-k] = \overline{\hat{f}[k]}
\end{equation}
\textit{Proof:} Conjugate the definition; $\overline{f} = f$ for real $f$; $\overline{e^{-ikt}} = e^{ikt} = e^{-i(-k)t}$. \quad
\textit{Intuition:} Negative frequencies are redundant for real signals. The power spectrum is symmetric.

\subsection{Scaling}
\begin{equation}
    \widehat{f(at)}[k] = \frac{1}{|a|} \hat{f}\!\left[\frac{k}{a}\right]
\end{equation}
\textit{Intuition:} Compressing in time stretches in frequency, and vice versa. Short pulses have broad spectra. This is the uncertainty principle.

% ======================================================================
\section{Gaussian Blur in the Frequency Domain}
% ======================================================================

The Gaussian kernel $g(x) = \frac{1}{\sqrt{2\pi}\sigma} e^{-x^2/(2\sigma^2)}$ has Fourier transform:
\begin{equation}
    \hat{g}[k] = e^{-2\pi^2 \sigma^2 k^2}
\end{equation}

By the convolution theorem, blurring $f$ with $g$ gives:
\begin{equation}
    \widehat{f * g}[k] = \hat{f}[k] \cdot e^{-2\pi^2 \sigma^2 k^2}
\end{equation}

Each frequency $k$ gets multiplied by $e^{-2\pi^2\sigma^2 k^2}$:
\begin{itemize}
    \item Small $k$ (low frequency): $e^{-\text{small}} \approx 1$ --- preserved
    \item Large $k$ (high frequency): $e^{-\text{large}} \approx 0$ --- destroyed
\end{itemize}

The \textbf{effective cutoff frequency} (where attenuation reaches $1/e$):
\begin{equation}
    k_{\text{cutoff}} = \frac{1}{2\pi\sigma}
\end{equation}

Larger $\sigma$ = lower cutoff = more aggressive filtering.

\begin{application}
This is the core mechanism of FB-CycleGAN. The steganographic generator hides information in high-$k$ coefficients (fine patterns invisible to the eye). The Gaussian blur multiplies those coefficients by $\approx 0$, destroying the hidden channel. The generator can only encode information in the coefficients that survive the blur, which have insufficient capacity to store the tumor coordinates needed for cycle reconstruction.
\end{application}

% ======================================================================
\section{Application: Steganographic Detection via FFT}
% ======================================================================

Given a generated image $\hat{h} = G(p)$ and its source $p$:

\begin{enumerate}
    \item Compute residual: $r = |\hat{h} - p|$
    \item Apply Hann window: $r_w = r \cdot w_{\mathrm{Hann}}$ (taper edges to zero to prevent spectral leakage---the DFT assumes periodic signals, and sharp boundaries at the image edges create artificial high-frequency content)
    \item Compute 2D DFT: $R = \mathrm{DFT}(r_w)$
    \item Compute power spectrum: $P = |R|^2$
    \item Shift zero-frequency to centre: $P_c = \mathrm{fftshift}(P)$
    \item Compute radial average $P(\rho)$ and plot on log-log scale
    \item Average over many samples for statistical significance
\end{enumerate}

\textbf{Steganographic generator:} Structured energy at high frequencies (bumps/plateaus in $P(\rho)$). The generator is encoding data as specific high-frequency patterns.

\textbf{Honest generator:} Steep rolloff at high frequencies, no structured peaks. The residual contains only natural image differences.

\textbf{Perturbation test:} Inject noise $\epsilon \sim \mathcal{N}(0, \sigma_\epsilon^2)$ before the reverse cycle: $\hat{p} = F(\hat{h} + \epsilon)$. A steganographic model collapses ($>2\times$ L1 increase) because the noise destroys the hidden signal. An honest model degrades gracefully.

\end{document}
